
\documentclass{article}
\usepackage{colortbl}
\usepackage{makecell}
\usepackage{multirow}
\usepackage{supertabular}

\begin{document}

\newcounter{utterance}

\centering \large Interaction Transcript for game `ifeval', experiment `ifeval', episode 161 with reflexion{-}llama{-}3.1{-}8b{-}instruct.
\vspace{24pt}

{ \footnotesize  \setcounter{utterance}{1}
\setlength{\tabcolsep}{0pt}
\begin{supertabular}{c@{$\;$}|p{.15\linewidth}@{}p{.15\linewidth}p{.15\linewidth}p{.15\linewidth}p{.15\linewidth}p{.15\linewidth}}
    \# & \multicolumn{2}{c}{Player} && \multicolumn{2}{c}{Game Master} \\
    \hline

    \theutterance \stepcounter{utterance}  
    & & & \multicolumn{4}{p{0.6\linewidth}}{
        \cellcolor[rgb]{0.9,0.9,0.9}{
            \makecell[{{p{\linewidth}}}]{
                \texttt{\tiny{[P1$\langle$GM]}}
                \texttt{You have attempted this same task before and did not succeed. These are the plans you wrote after those attempts.} \\
\texttt{Trial \#1: New plan: } \\
\\ 
\texttt{Upon reflection, I realize that my previous attempt was too focused on describing the performance itself, without considering the specific task of continuing and expanding the given sentence. I got caught up in the excitement of the performance and lost sight of the goal.} \\
\\ 
\texttt{To improve, I will take a step back and re{-}examine the original sentence: "A man is performing for a crowd on a stage, ..." My task is to continue and expand this sentence, not to create a new story or describe a different scene.} \\
\\ 
\texttt{Next time, I will focus on building on the existing sentence, using the given information as a starting point. I will ask myself questions like: What is the man doing on stage? What is the crowd's reaction? What is the atmosphere like?} \\
\\ 
\texttt{To achieve this, I will take a more structured approach, breaking down the task into smaller, manageable chunks. I will start by identifying the key elements of the original sentence, such as the man, the crowd, and the stage. Then, I will use these elements as a foundation to build on, adding new details and descriptions to continue and expand the sentence.} \\
\\ 
\texttt{For example, I might start by describing the man's actions on stage, such as singing or playing an instrument. I might then describe the crowd's reaction, including their emotions and physical responses. Finally, I might add more details about the stage and the atmosphere, using sensory language to bring the scene to life.} \\
\\ 
\texttt{By taking a more focused and structured approach, I hope to create a more cohesive and engaging continuation of the original sentence.} \\
\\ 
\texttt{Trial \#2: New plan: } \\
\\ 
\texttt{Upon reflection, I realize that my previous attempt was too focused on describing the performance itself, without considering the specific task of continuing and expanding the given sentence. I got caught up in the excitement of the performance and lost sight of the goal.} \\
\\ 
\texttt{To improve, I will take a step back and re{-}examine the original sentence: "A man is performing for a crowd on a stage, ..." My task is to continue and expand this sentence, not to create a new story or describe a different scene.} \\
\\ 
\texttt{Next time, I will focus on building on the existing sentence, using the given information as a starting point. I will ask myself questions like: What is the man doing on stage? What is the crowd's reaction? What is the atmosphere like?} \\
\\ 
\texttt{To achieve this, I will take a more structured approach, breaking down the task into smaller, manageable chunks. I will start by identifying the key elements of the original sentence, such as the man, the crowd, and the stage. Then, I will use these elements as a foundation to build on, adding new details and descriptions to continue and expand the sentence.} \\
\\ 
\texttt{However, I also realize that my previous attempt was too repetitive and lacked variety. I repeated the same descriptions and ideas multiple times, without adding any new insights or perspectives.} \\
\\ 
\texttt{To address this, I will make a conscious effort to vary my descriptions and add more depth and nuance to the scene. I will use sensory language to bring the scene to life, incorporating details about the sights, sounds, and emotions of the performance.} \\
\\ 
\texttt{For example, I might describe the man's voice as "rich and velvety," or the crowd's reaction as "electric and frenzied." I might also add more details about the stage and the atmosphere, such as the smell of sweat and smoke, or the feeling of the crowd's energy.} \\
\\ 
\texttt{By taking a more focused and structured approach, and varying my descriptions to add more depth and nuance, I hope to create a more engaging and immersive continuation of the original sentence.} \\
\\ 
\texttt{Trial \#3: New plan: } \\
\\ 
\texttt{Upon reflection, I realize that my previous attempts were too focused on describing the performance itself, without considering the specific task of continuing and expanding the given sentence. I got caught up in the excitement of the performance and lost sight of the goal.} \\
\\ 
\texttt{To improve, I will take a step back and re{-}examine the original sentence: "A man is performing for a crowd on a stage, ..." My task is to continue and expand this sentence, not to create a new story or describe a different scene.} \\
\\ 
\texttt{Next time, I will focus on building on the existing sentence, using the given information as a starting point. I will ask myself questions like: What is the man doing on stage? What is the crowd's reaction? What is the atmosphere like?} \\
\\ 
\texttt{To achieve this, I will take a more structured approach, breaking down the task into smaller, manageable chunks. I will start by identifying the key elements of the original sentence, such as the man, the crowd, and the stage. Then, I will use these elements as a foundation to build on, adding new details and descriptions to continue and expand the sentence.} \\
\\ 
\texttt{However, I also realize that my previous attempts were too repetitive and lacked variety. I repeated the same descriptions and ideas multiple times, without adding any new insights or perspectives.} \\
\\ 
\texttt{To address this, I will make a conscious effort to vary my descriptions and add more depth and nuance to the scene. I will use sensory language to bring the scene to life, incorporating details about the sights, sounds, and emotions of the performance.} \\
\\ 
\texttt{Additionally, I will pay closer attention to the sentence structure and grammar, making sure to use proper punctuation and sentence flow to create a cohesive and engaging narrative.} \\
\\ 
\texttt{To achieve this, I will take a more focused and structured approach, varying my descriptions to add more depth and nuance, and paying closer attention to sentence structure and grammar. I will also make a conscious effort to add more variety and creativity to my descriptions, using sensory language and vivid imagery to bring the scene to life.} \\
\\ 
\texttt{For example, I might start by describing the man's actions on stage, such as singing or playing an instrument. I might then describe the crowd's reaction, including their emotions and physical responses. Finally, I might add more details about the stage and the atmosphere, using sensory language to bring the scene to life.} \\
\\ 
\texttt{By taking a more focused and structured approach, varying my descriptions to add more depth and nuance, and paying closer attention to sentence structure and grammar, I hope to create a more engaging and immersive continuation of the original sentence.} \\
\\ 
\\ 
\texttt{Can you please continue and expand: "A man is performing for a crowd on a stage, ..."} \\
\texttt{Use at least 60 sentences, but less than 600 words.} \\
            }
        }
    }
     \\ \\

    \theutterance \stepcounter{utterance}  
    & \multicolumn{4}{p{0.6\linewidth}}{
        \cellcolor[rgb]{0.561,0.737,0.561}{
            \makecell[{{p{\linewidth}}}]{
                \texttt{\tiny{[P1$\rangle$GM]}}
                \texttt{A man is performing for a crowd on a stage, his voice echoing through the microphone as he sings with passion and conviction. The crowd is entranced, their faces aglow with excitement and anticipation. They sway to the rhythm, their hands clapping in time with the beat.} \\
\\ 
\texttt{The man's voice is rich and velvety, his words pouring out like honey as he weaves a spell of enchantment over the audience. His eyes scan the crowd, locking onto individual faces and holding their gaze for a moment before moving on. He is a master of his craft, a true performer who knows how to work a crowd.} \\
\\ 
\texttt{The stage is dimly lit, the only light coming from a single spotlight that shines down on the man. The air is thick with the smell of sweat and smoke, the crowd's energy palpable as they cheer and whistle for more. The man's music is a fusion of rock and blues, a unique sound that is both familiar and yet completely new.} \\
\\ 
\texttt{As he sings, the man's body moves in time with the music, his hips swaying and his arms waving in the air. He is a force of nature, a whirlwind of energy and emotion that sweeps the crowd up in its wake. The crowd is caught up in the moment, their worries and cares forgotten as they lose themselves in the music.} \\
\\ 
\texttt{The man's voice soars, reaching heights that are almost impossible to believe. He is a true virtuoso, a master of his craft who has spent years honing his skills. The crowd is in awe, their eyes wide with wonder as they watch him perform.} \\
\\ 
\texttt{As the song comes to an end, the crowd erupts into cheers and applause, their hands clapping and whistling for more. The man takes a bow, his face flushed with excitement and his eyes shining with pride. He knows he has given it his all, pouring his heart and soul into the performance.} \\
\\ 
\texttt{The crowd is on its feet, cheering and whistling as the man takes his final bow. The stage is bathed in a warm, golden light, the only sound the thunderous applause of the crowd. The man is a true rock star, a legend in his own right who has given the performance of a lifetime.} \\
\\ 
\texttt{As the applause dies down, the man takes a moment to catch his breath, his chest heaving with exertion. He looks out at the crowd, his eyes scanning the sea of faces as he searches for someone, anyone, who understands him. And then, in a flash of inspiration, he knows exactly what to do.} \\
\\ 
\texttt{He launches into a new song, one that is even more energetic and passionate than the last. The crowd is caught off guard, their eyes wide with surprise as they watch the man perform. They are in awe, their faces aglow with excitement and anticipation.} \\
\\ 
\texttt{The man's voice is a force of nature, a whirlwind of energy and emotion that sweeps the crowd up in its wake. He is a true performer, a master of his craft who knows how to work a crowd. The crowd is entranced, their faces aglow with excitement and anticipation.} \\
\\ 
\texttt{As the song comes to an end, the crowd erupts into cheers and applause, their hands clapping and whistling for more. The man takes a bow, his face flushed with excitement and his eyes shining with pride. He knows he has given it his all, pouring his heart and soul into the performance.} \\
\\ 
\texttt{The crowd is on its feet, cheering and whistling as the man takes his final bow. The stage is bathed in a warm, golden light, the only sound the thunderous applause of the crowd. The man is a true rock star, a legend in his own right who has given the performance of a lifetime.} \\
\\ 
\texttt{As the applause dies down, the man takes a moment to catch his breath, his chest heaving with exertion. He looks out at the crowd, his eyes scanning the sea of faces as he searches for someone, anyone, who understands him. And then, in a flash of inspiration, he knows exactly what to do.} \\
\\ 
\texttt{He launches into a new song, one that is even more energetic and passionate than the last. The crowd is caught off guard, their eyes wide with surprise as they watch the man perform. They are in awe, their faces aglow with excitement and anticipation.} \\
\\ 
\texttt{The man's voice is a force of nature, a whirlwind of energy and emotion that sweeps the crowd up in its wake. He is a true performer, a master of his craft who knows how to work a crowd. The crowd is entranced, their faces aglow with excitement and anticipation.} \\
\\ 
\texttt{As the song comes to an end, the crowd erupts into cheers and applause, their hands clapping and whistling for more. The man takes a bow, his face flushed with excitement and his eyes shining with pride. He knows he has given it his all, pouring his heart and soul into the performance.} \\
\\ 
\texttt{The crowd is on its feet, cheering and whistling as the man takes his final bow. The stage is bathed in a warm, golden light, the only sound the thunderous applause of the crowd. The man is a true rock star, a legend in his own right who has given the performance of a lifetime.} \\
\\ 
\texttt{As the applause dies down, the man takes a moment to catch his breath, his chest heaving with exertion. He looks out at the crowd, his eyes scanning the sea of faces as he searches for someone, anyone, who understands him. And then, in a flash of inspiration, he knows exactly what to do.} \\
\\ 
\texttt{He launches into a new song, one that is even more energetic and passionate than the last. The crowd is caught off guard, their eyes wide with surprise as they watch the man perform. They are in awe, their faces aglow with excitement and anticipation.} \\
\\ 
\texttt{The man's voice is a force of nature, a whirlwind of energy and emotion that sweeps the crowd up in its wake. He is a true performer, a master of his craft who knows how to work a crowd. The crowd is entranced, their faces aglow with excitement and anticipation.} \\
\\ 
\texttt{As the song comes to an end, the crowd erupts into cheers and applause, their hands clapping and whistling for more. The man takes} \\
            }
        }
    }
    & & \\ \\

    \theutterance \stepcounter{utterance}  
    & & & \multicolumn{2}{p{0.3\linewidth}}{
        \cellcolor[rgb]{0.9,0.9,0.9}{
            \makecell[{{p{\linewidth}}}]{
                \texttt{\tiny{[GM$|$GM]}}
                {'length_constraints:number_sentences': {'relation': 'at least', 'num_sentences': 60}, 'length_constraints:number_words': {'relation': 'less than', 'num_words': 600}}
            }
        }
    }
    & & \\ \\

    \theutterance \stepcounter{utterance}  
    & & & \multicolumn{2}{p{0.3\linewidth}}{
        \cellcolor[rgb]{0.9,0.9,0.9}{
            \makecell[{{p{\linewidth}}}]{
                \texttt{\tiny{[GM$|$GM]}}
                \texttt{game\_result = LOSE} \\
            }
        }
    }
    & & \\ \\

    \theutterance \stepcounter{utterance}  
    & & \multicolumn{4}{p{0.6\linewidth}}{
        \cellcolor[rgb]{0.95,0.95,0.95}{
            \makecell[{{p{\linewidth}}}]{
                \texttt{\tiny{[P1$\leftrightarrow$P]}}
                \texttt{Post{-}episode Reflexion (not sent to GM):} \\
\texttt{New plan: } \\
\\ 
\texttt{Upon reflection, I realize that my previous attempts were too focused on describing the performance itself, without considering the specific task of continuing and expanding the given sentence. I got caught up in the excitement of the performance and lost sight of the goal.} \\
\\ 
\texttt{To improve, I will take a step back and re{-}examine the original sentence: "A man is performing for a crowd on a stage, ..." My task is to continue and expand this sentence, not to create a new story or describe a different scene.} \\
\\ 
\texttt{Next time, I will focus on building on the existing sentence, using the given information as a starting point. I will ask myself questions like: What is the man doing on stage? What is the crowd's reaction? What is the atmosphere like?} \\
\\ 
\texttt{To achieve this, I will take a more structured approach, breaking down the task into smaller, manageable chunks. I will start by identifying the key elements of the original sentence, such as the man, the crowd, and the stage. Then, I will use these elements as a foundation to build on, adding new details and descriptions to continue and expand the sentence.} \\
\\ 
\texttt{However, I also realize that my previous attempts were too repetitive and lacked variety. I repeated the same descriptions and ideas multiple times, without adding any new insights or perspectives.} \\
\\ 
\texttt{To address this, I will make a conscious effort to vary my descriptions and add more depth and nuance to the scene. I will use sensory language to bring the scene to life, incorporating details about the sights, sounds, and emotions of the performance.} \\
\\ 
\texttt{Additionally, I will pay closer attention to the sentence structure and grammar, making sure to use proper punctuation and sentence flow to create a cohesive and engaging narrative.} \\
\\ 
\texttt{To achieve this, I will take a more focused and structured approach, varying my descriptions to add more depth and nuance, and paying closer attention to sentence structure and grammar. I will also make a conscious effort to add more variety and creativity to my descriptions, using sensory language and vivid imagery to bring the scene to life.} \\
\\ 
\texttt{To specifically address the issue of repetition, I will make a conscious effort to introduce new ideas and perspectives in each iteration, rather than relying on the same descriptions and ideas. I will also pay closer attention to the sentence structure and grammar, making sure to use proper punctuation and sentence flow to create a cohesive and engaging narrative.} \\
\\ 
\texttt{For example, I might start by describing the man's actions on stage, such as singing or playing an instrument. I might then describe the crowd's reaction, including their emotions and physical responses. Finally, I might add more details about the stage and the atmosphere, using sensory language to bring the scene to life.} \\
\\ 
\texttt{By taking a more focused and structured approach, varying my descriptions to add more depth and nuance, and paying closer attention to sentence structure and grammar, I hope to create a more engaging and immersive continuation of the original sentence.} \\
            }
        }
    }
    & \\ \\

\end{supertabular}
}

\end{document}
